\documentclass[conference]{IEEEtran}
\usepackage{cite}
\usepackage{amsmath,amssymb,amsfonts}
\usepackage{algorithmic}
\usepackage{graphicx}
\usepackage{textcomp}
\usepackage{xcolor}
\usepackage{hyperref}

\def\BibTeX{{\rm B\kern-.05em{\sc i\kern-.025em b}\kern-.08em
    T\kern-.1667em\lower.7ex\hbox{E}\kern-.125emX}}
\begin{document}

\title{AttenX: Enhancing Fine-Grained Text-to-Image Synthesis with Self-Attention and Spectral Normalization}

\author{\IEEEauthorblockN{Ahmed Islam Farouk}
\IEEEauthorblockA{\textit{Department of Computer Science} \\
\textit{Lab Module 6} \\
ahmedislamfarouk@example.com}
}

\maketitle

\begin{abstract}
Generating high-fidelity images from natural language descriptions is a significant challenge in computer vision. While the Attentional Generative Adversarial Network (AttnGAN) introduced fine-grained attention mechanisms, it often suffers from a lack of global structural coherence and training instability. This paper proposes AttenX, an enhanced architecture that integrates Self-Attention modules into the generator to capture long-range dependencies and applies Spectral Normalization to the multi-scale discriminators to stabilize training. Experimental results on the CUB-200-2011 dataset demonstrate that AttenX achieves a 12.16\% improvement in Inception Score (IS) and a 22.64\% reduction in Fréchet Inception Distance (FID) compared to the baseline AttnGAN.
\end{abstract}

\begin{IEEEkeywords}
Generative Adversarial Networks, Text-to-Image Synthesis, Self-Attention, Spectral Normalization, AttnGAN.
\end{IEEEkeywords}

\section{Introduction}
Text-to-Image synthesis involves generating a visual representation that aligns with a given natural language prompt. Recent advancements, such as StackGAN and AttnGAN, have utilized multi-stage architectures to progressively refine image resolution. AttnGAN, in particular, introduced a word-level attention mechanism that allows the generator to focus on specific words when synthesizing regions of an image.

However, AttnGAN has two primary drawbacks: (1) it lacks a mechanism to maintain global structural coherence across the entire image, often leading to anatomical distortions in complex objects (e.g., birds with multiple beaks), and (2) the discriminators often become too powerful early in training, causing mode collapse or vanishing gradients.

In this paper, we present AttenX, which addresses these issues by incorporating non-local self-attention and spectral normalization. We demonstrate through quantitative and qualitative metrics that these enhancements significantly improve the realism and stability of the synthesis process.

\section{Proposed Methodology}
The AttenX architecture is built upon the foundational AttnGAN framework but introduces three key modifications.

\subsection{Self-Attention Integration}
To capture global dependencies, we inject a Self-Attention module after the $128 \times 128$ resolution stage. This stage was selected to balance computational efficiency with structural refinement. The attention mechanism computes a context vector by considering the relationship between every pair of pixels in the feature map, allowing the model to ensure that distant regions (e.g., the head and tail of a bird) are anatomically consistent.

The mathematical formulation follows the non-local neural network approach:
\begin{equation}
Y = \gamma \cdot \text{Softmax}(Q^T K) V + X
\end{equation}
where $Q, K, V$ are linear projections of the input $X$, and $\gamma$ is a learnable scale parameter initialized to zero.

\subsection{Spectral Normalization in Discriminators}
To stabilize training, we apply Spectral Normalization (SN) to all convolutional layers in the three multi-scale discriminators ($D_0, D_1, D_2$). SN constrains the Lipschitz constant of the discriminator by normalizing the spectral norm of the weight matrices, preventing the gradients from exploding and ensuring a smoother loss landscape.

\subsection{Two-Time Scale Update Rule (TTUR)}
With the stabilization provided by SN, we employ TTUR, using a higher learning rate for the discriminators ($4 \times 10^{-4}$) than for the generator ($1 \times 10^{-4}$). This allows the discriminators to provide more informative gradients to the generator throughout the training process.

\section{Experimental Results}
We evaluated AttenX on the CUB-200-2011 bird dataset. 

\subsection{Quantitative Analysis}
We compared AttenX against the original AttnGAN baseline using the Inception Score (IS) and Fréchet Inception Distance (FID). As shown in Table \ref{tab:metrics}, AttenX significantly outperforms the baseline.

\begin{table}[htbp]
\caption{Quantitative Performance Comparison}
\begin{center}
\begin{tabular}{|l|c|c|}
\hline
\textbf{Model} & \textbf{IS ($\uparrow$)} & \textbf{FID ($\downarrow$)} \\
\hline
AttnGAN (Baseline) & 4.36 & 23.54 \\
\hline
\textbf{AttenX (Ours)} & \textbf{4.89} & \textbf{18.21} \\
\hline
\end{tabular}
\label{tab:metrics}
\end{center}
\end{table}

\subsection{Qualitative Analysis}
Visual inspection of the generated samples reveals that AttenX produces images with superior global coherence. Anatomical inconsistencies common in the baseline, such as misplaced limbs or distorted beaks, are significantly reduced in the AttenX outputs.

\section{Conclusion}
This study presented AttenX, an enhanced version of the AttnGAN model for text-to-image synthesis. By integrating self-attention and spectral normalization, we effectively addressed the issues of structural incoherence and training instability. The quantitative results demonstrate a clear statistical improvement in image quality and distribution alignment. Future work will explore the integration of transformer-based encoders for even finer text-image alignment.

\section*{References}
\begin{enumerate}
\item T. Xu et al., "AttnGAN: Fine-Grained Text to Image Generation with Attentional Generative Adversarial Networks," \textit{Proc. CVPR}, 2018. DOI: \href{https://doi.org/10.1109/CVPR.2018.00143}{10.1109/CVPR.2018.00143}
\item H. Zhang et al., "Self-Attention Generative Adversarial Networks," \textit{Proc. ICML}, 2019. DOI: \href{https://doi.org/10.48550/arXiv.1805.08318}{10.48550/arXiv.1805.08318}
\item T. Miyato et al., "Spectral Normalization for Generative Adversarial Networks," \textit{Proc. ICLR}, 2018. DOI: \href{https://doi.org/10.48550/arXiv.1802.05957}{10.48550/arXiv.1802.05957}
\item M. Heusel et al., "GANs Trained by a Two Time-Scale Update Rule Converge to a Local Nash Equilibrium," \textit{Proc. NIPS}, 2017. DOI: \href{https://doi.org/10.48550/arXiv.1706.08500}{10.48550/arXiv.1706.08500}
\item H. Zhang et al., "StackGAN++: Realistic Image Synthesis with Stacked Generative Adversarial Networks," \textit{IEEE Trans. PAMI}, 2018. DOI: \href{https://doi.org/10.1109/TPAMI.2018.2856256}{10.1109/TPAMI.2018.2856256}
\end{enumerate}

\end{document}
